\documentclass[12pt]{ximera}
\title{Practice Exam 2}
\begin{document}
\begin{abstract}
Practice for Exam 2, covering Chapter 2: critical players and veto power, counting
coalitions, Banzhaf power and dummies, reconstructing winning coalitions from pivots,
and Shapley-Shubik power in a law firm.
\end{abstract}
\maketitle

\section*{Practice Exam 2}

\begin{problem}
Assess whether each of the following claims is true or false.

\begin{prompt}
\textbf{(a)} Claim: in a weighted voting system, it is possible for every player in a
\emph{winning} coalition to be critical.

\begin{multipleChoice}
\choice[correct]{True}
\choice{False}
\end{multipleChoice}

\begin{feedback}
True. In a winning coalition that only just reaches the quota, losing any member drops
it below. For example, in $[12:8,6,4,2]$ the coalition $\set{B,C,D}$ has weight exactly
$12$, so all three members are critical.
\end{feedback}
\end{prompt}

\begin{prompt}
\textbf{(b)} Claim: in a weighted voting system, it is possible for a player with weight
less than half of the quota to have veto power.

\begin{multipleChoice}
\choice[correct]{True}
\choice{False}
\end{multipleChoice}

\begin{feedback}
True. Veto power only requires that the other players together fall short of the quota.
In $[10:4,3,2,1]$ the quota equals the total weight, so every player --- including the
one with weight $1$, far below half of $10$ --- can block any motion.
\end{feedback}
\end{prompt}

\begin{prompt}
\textbf{(c)} Claim: in a weighted voting system, if two players have the same number of
votes, then they have equal power in the system.

\begin{multipleChoice}
\choice[correct]{True}
\choice{False}
\end{multipleChoice}

\begin{feedback}
True. Two players with equal weight are interchangeable: swapping them turns every
winning coalition into another winning coalition, so each is critical (or pivotal)
exactly as often as the other.
\end{feedback}
\end{prompt}
\end{problem}

\begin{problem}
Use counting of subsets and sequences to answer the following.

\begin{prompt}
\textbf{(a)} In a system with $8$ players, how many coalitions have two or more
players?

$\answer{247}$

\begin{feedback}
$2^8=256$ subsets, minus the empty one and the $8$ single-player ones:
$256-8-1=247$.
\end{feedback}
\end{prompt}

\begin{prompt}
\textbf{(b)} In a system with $11$ players, how many coalitions include $P_1$, $P_2$,
and $P_3$ but exclude $P_{10}$ and $P_{11}$?

$\answer{64}$

\begin{feedback}
Five players are settled, leaving $6$ free: $2^6=64$.
\end{feedback}
\end{prompt}

\begin{prompt}
\textbf{(c)} In the system $[28:7,6,5,4,3,2,1]$ with seven players, how many
sequential coalitions are there?

$\answer{5040}$

\begin{feedback}
$7!=5040$.
\end{feedback}
\end{prompt}
\end{problem}

\begin{problem}
Consider the weighted voting system $[10:7,5,3,1]$ with players A, B, C, and D.

\begin{prompt}
\textbf{(a)} How many winning coalitions are there?

$\answer{6}$

\begin{feedback}
The winning coalitions and their critical players:
\[
\begin{array}{ll}
\set{A,B}=12 & \text{A, B critical}\\
\set{A,C}=10 & \text{A, C critical}\\
\set{A,B,C}=15 & \text{A critical}\\
\set{A,B,D}=13 & \text{A, B critical}\\
\set{A,C,D}=11 & \text{A, C critical}\\
\set{A,B,C,D}=16 & \text{A critical}
\end{array}
\]
\end{feedback}
\end{prompt}

\begin{prompt}
\textbf{(b)} Determine the Banzhaf power distribution, as percentages.

A: $\answer{60}\%$ \quad B: $\answer{20}\%$ \quad C: $\answer{20}\%$ \quad D: $\answer{0}\%$

\begin{feedback}
Critical counts: A $6$, B $2$, C $2$, D $0$, total $10$.
\end{feedback}
\end{prompt}

\begin{prompt}
\textbf{(c)} Which players, if any, are dummies?

\begin{selectAll}
\choice{A}
\choice{B}
\choice{C}
\choice[correct]{D}
\end{selectAll}

\begin{feedback}
D is never critical --- every winning coalition containing D still wins without it ---
so D is a dummy, despite having a vote.
\end{feedback}
\end{prompt}

\begin{prompt}
\textbf{(d)} Claim: this system is \emph{equivalent} to $[3:2,2,1,0]$.

\begin{multipleChoice}
\choice{True}
\choice[correct]{False}
\end{multipleChoice}

\begin{feedback}
False. In $[3:2,2,1,0]$ the coalition $\set{B,C}$ has weight $3$ and wins, but in
$[10:7,5,3,1]$ it has weight $8$ and loses.

(An equivalent system does exist: $[4:3,2,1,0]$ has exactly the same six winning
coalitions.)
\end{feedback}
\end{prompt}
\end{problem}

\begin{problem}
The following are all $4!=24$ sequential coalitions for a weighted voting system with
four players A, B, C, and D. In some of them the pivotal player is underlined.
\[
\begin{array}{cccc}
\seq{A,\underline{B},C,D} & \seq{B,A,C,D} & \seq{C,A,B,D} & \seq{D,A,B,C}\\
\seq{A,B,D,C} & \seq{B,A,D,C} & \seq{C,A,D,B} & \seq{D,A,C,B}\\
\seq{A,C,\underline{B},D} & \seq{B,C,\underline{A},D} & \seq{C,B,A,D} & \seq{D,B,A,C}\\
\seq{A,C,\underline{D},B} & \seq{B,C,D,\underline{A}} & \seq{C,B,D,A} & \seq{D,B,C,A}\\
\seq{A,D,\underline{B},C} & \seq{B,D,\underline{A},C} & \seq{C,D,\underline{A},B} & \seq{D,C,A,B}\\
\seq{A,D,\underline{C},B} & \seq{B,D,C,A} & \seq{C,D,B,A} & \seq{D,C,B,A}
\end{array}
\]

\begin{prompt}
\textbf{(a)} Based on the marked pivots, which are the five winning coalitions?

\begin{selectAll}
\choice[correct]{$\set{A,B}$}
\choice{$\set{A,C}$}
\choice{$\set{A,D}$}
\choice{$\set{B,C,D}$}
\choice[correct]{$\set{A,B,C}$}
\choice[correct]{$\set{A,B,D}$}
\choice[correct]{$\set{A,C,D}$}
\choice[correct]{$\set{A,B,C,D}$}
\end{selectAll}

\begin{hint}
A pivot in position $k$ means the first $k$ players win, and the first $k-1$ lose.
\end{hint}

\begin{feedback}
$\seq{A,\underline{B},\ldots}$ gives $\set{A,B}$ winning. $\seq{A,C,\underline{B},D}$ gives
$\set{A,C}$ losing and $\set{A,B,C}$ winning. $\seq{A,C,\underline{D},B}$ gives
$\set{A,C,D}$ winning. $\seq{A,D,\underline{B},C}$ gives $\set{A,D}$ losing and
$\set{A,B,D}$ winning. $\seq{B,C,D,\underline{A}}$ gives $\set{B,C,D}$ losing. With the
grand coalition, that is five winning coalitions.
\end{feedback}
\end{prompt}

\begin{prompt}
\textbf{(b)} Now find the pivotal player in every sequential coalition. How many times
is each player pivotal?

A: $\answer{14}$ \quad B: $\answer{6}$ \quad C: $\answer{2}$ \quad D: $\answer{2}$

\begin{feedback}
Every winning coalition contains A, so the pivot is never earlier than A. A itself is
the pivot exactly when the players ahead of it already complete a winning coalition with
A --- that is, when B comes before A, or when both C and D do. That happens in $14$
orderings, e.g.\ $\seq{B,\underline{A},C,D}$ or $\seq{C,D,\underline{A},B}$.

In the other $10$ orderings, the pivot is the player who completes $\set{A,B}$ or
$\set{A,C,D}$ after A:
\begin{itemize}
\item B is pivotal in $\seq{A,B,C,D}$, $\seq{A,B,D,C}$, $\seq{A,C,B,D}$, $\seq{A,D,B,C}$,
      $\seq{C,A,B,D}$, $\seq{D,A,B,C}$ --- $6$ times;
\item C is pivotal in $\seq{A,D,C,B}$ and $\seq{D,A,C,B}$ --- $2$ times;
\item D is pivotal in $\seq{A,C,D,B}$ and $\seq{C,A,D,B}$ --- $2$ times.
\end{itemize}
The counts total $14+6+2+2=24$. \checkmark
\end{feedback}
\end{prompt}

\begin{prompt}
\textbf{(c)} Determine the Shapley-Shubik power distribution, as percentages rounded to
two decimal places.

A: $\answer[tolerance=0.005]{58.33}\%$ \quad B: $\answer[tolerance=0.005]{25}\%$ \quad
C: $\answer[tolerance=0.005]{8.33}\%$ \quad D: $\answer[tolerance=0.005]{8.33}\%$

\begin{feedback}
$\frac{14}{24}\approx 58.33\%$, $\frac{6}{24}=25\%$, $\frac{2}{24}\approx 8.33\%$ each for C
and D. One weighted system with these winning coalitions is $[5:3,2,1,1]$.
\end{feedback}
\end{prompt}
\end{problem}

\begin{problem}
Consider a law firm run by four partners: $A$, the senior partner, and junior partners
$B_1$, $B_2$, $B_3$. Each partner has one vote, and decisions can be made by a simple
majority or by the senior partner with any one junior partner's support. Use $B$ for
any one junior partner.

\begin{prompt}
\textbf{(a)} What is the total number of sequential coalitions in this four-player
system?

$\answer{24}$

\begin{feedback}
$4!=24$.
\end{feedback}
\end{prompt}

\begin{prompt}
\textbf{(b)} There are four sequential coalition types, depending on where A falls. In
how many is A pivotal?

$\answer{2}$

\begin{feedback}
\[
\begin{array}{ll}
\seq{A,\underline{B},B,B} & \text{A with one junior wins}\\
\seq{B,\underline{A},B,B} & \text{same}\\
\seq{B,B,\underline{A},B} & \text{two juniors are not a majority; A completes it}\\
\seq{B,B,\underline{B},A} & \text{three of four is a majority}
\end{array}
\]
\end{feedback}
\end{prompt}

\begin{prompt}
\textbf{(c)} How many distinct sequential coalitions are there of each type --- that is,
how many ways are there to arrange the junior partners?

$\answer{6}$

\begin{feedback}
The three juniors can be ordered $3!=6$ ways within a type, and $4\cdot 6=24=4!$.
\end{feedback}
\end{prompt}

\begin{prompt}
\textbf{(d)} Determine the Shapley-Shubik power distribution, as percentages rounded to
two decimal places.

A: $\answer[tolerance=0.005]{50}\%$ \quad Each junior partner:
$\answer[tolerance=0.005]{16.67}\%$

\begin{feedback}
A is pivotal in $2\cdot 6=12$ of $24$: $50\%$. The juniors share the other $12$, so $4$
each: $\frac{4}{24}\approx 16.67\%$.
\end{feedback}
\end{prompt}
\end{problem}

\end{document}
